\begin{tikzpicture}[x=1cm,y=0.75cm,font=\small]
\node[seqactor] (sensor) at (0,0) {Vlákno senzora};
\node[seqactor] (server) at (9,0) {Server: vynechať 5 ACK};
\draw[dashed] (sensor.south)--(0,-15.7);
\draw[dashed] (server.south)--(9,-15.7);
\foreach \y/\label in {-1/{0 s: DATA seq=2, pokus 1},-2.5/{3 s: DATA seq=2, pokus 2},-4/{6 s: DATA seq=2, pokus 3}} {
 \draw[seqarrow] (0,\y)--node[above]{\label}(9,\y);
}
\node[seqnote,text width=6.4cm] at (4.5,-5.6) {ACK 1--3 potlačené.\\9 s: zahodiť čakanie na seq=2.};
\foreach \y/\label in {-7.2/{10 s: DATA seq=3, pokus 1},-8.7/{13 s: DATA seq=3, pokus 2},-10.2/{16 s: DATA seq=3, pokus 3}} {
 \draw[seqarrow] (0,\y)--node[above]{\label}(9,\y);
}
\node[seqnote,text width=6cm] at (4.5,-11.6) {ACK 4--5 potlačené; šiesty pokus sa potvrdí.};
\draw[seqarrow] (9,-13.5)--node[above]{ACK seq=3}(0,-13.5);
\node[align=center,text width=8cm] at (4.5,-14.7) {Každé opakovanie zachová pôvodné bajty DATA.\\Ostatné senzory bežia nezávisle.};
\end{tikzpicture}
